\documentclass[11pt]{article}
\usepackage[margin=1in]{geometry}
\usepackage{amsmath,amssymb,physics,bm}
\usepackage[hidelinks]{hyperref}
\usepackage{microtype}
\newcommand{\epsz}{\epsilon_0}
\newcommand{\muz}{\mu_0}
\begin{document}
\section{Dispersion}\label{sec:em_fundementals_dispersion}
\textbf{Dispersion} is the phenomenon in which waves of different frequencies propagate at different speeds. More precisely, the relation between angular frequency $\omega$ and wavenumber $k$ is called the \emph{dispersion relation}, $\omega=\omega(k)$. A wave packet propagates at the group velocity $v_g=\dv{\omega}{k}$, while individual phase fronts move at $v_\mathrm{ph}=\omega/k$.

\subsection{Dispersion in free space}
In vacuum, where there are no charges or currents, Maxwell's equations give
\begin{align}
\nabla\cdot\mathbf E&=0,&\nabla\cdot\mathbf B&=0,\\
\nabla\times\mathbf E&=-\pdv{\mathbf B}{t},&
\nabla\times\mathbf B&=\muz\epsz\pdv{\mathbf E}{t}.
\end{align}
Take the curl of Faraday's law:
\begin{equation}
\nabla\times(\nabla\times\mathbf E)=-\pdv{}{t}(\nabla\times\mathbf B).
\end{equation}
Using $\nabla\times(\nabla\times\mathbf E)=\nabla(\nabla\cdot\mathbf E)-\nabla^2\mathbf E$, and substituting Amp\`ere--Maxwell's law, we obtain
\begin{equation}
\boxed{\nabla^2\mathbf E=\muz\epsz\pdv[2]{\mathbf E}{t}.}
\end{equation}
Try a plane wave $\mathbf E=\mathbf E_0 e^{i(\mathbf k\cdot\mathbf r-\omega t)}$. Then $\nabla^2\mathbf E=-k^2\mathbf E$ and $\partial_t^2\mathbf E=-\omega^2\mathbf E$, so
\begin{equation}
k^2=\muz\epsz\omega^2,\qquad
\boxed{\omega=ck},\qquad c=\frac{1}{\sqrt{\muz\epsz}}.
\end{equation}
Therefore
\begin{equation}
\boxed{v_\mathrm{ph}=\frac{\omega}{k}=c,\qquad v_g=\dv{\omega}{k}=c.}
\end{equation}
All frequencies propagate at the same speed: vacuum is nondispersive.


\subsection{The plasma frequency: physical intuition}

Imagine a uniform slab of plasma, with electron number density $n_e$
and stationary positive ions. Displace the electron fluid a small
distance $x$ to the right relative to the ions. The slab remains
neutral in its interior, but thin layers of opposite surface charge
appear at its edges, with magnitudes $\sigma_s=n_e e x$
(charge per unit area).

\begin{figure}[htbp]
\centering
\begin{tikzpicture}[>=Latex,scale=1.0]

% Original ion slab
\fill[red!7] (0,0) rectangle (6,2.6);
\draw[red!60!black,thick] (0,0) rectangle (6,2.6);

% Displaced electron slab
\fill[blue!12,opacity=0.65] (0.65,0) rectangle (6.65,2.6);
\draw[blue!70!black,thick,dashed]
    (0.65,0) rectangle (6.65,2.6);

% Positive ions
\foreach \xx in {0.3,1.1,...,5.9} {
  \foreach \yy in {0.35,1.05,1.75,2.35} {
    \node[red!75!black,font=\small] at (\xx,\yy) {$+$};
  }
}

% Electrons
\foreach \xx in {0.95,1.75,...,6.55} {
  \foreach \yy in {0.55,1.25,1.95,2.45} {
    \node[blue!80!black,font=\small] at (\xx,\yy) {$-$};
  }
}

% Surface charge regions
\fill[red!25,opacity=0.6] (0,0) rectangle (0.65,2.6);
\fill[blue!25,opacity=0.6] (6,0) rectangle (6.65,2.6);

% Sheet labels
\node[red!70!black,above] at (0.32,2.7)
    {$+\sigma_s$};
\node[blue!70!black,above] at (6.32,2.7)
    {$-\sigma_s$};

% Displacement arrow
\draw[->,very thick,blue!70!black]
    (2.7,3.3) -- (3.35,3.3)
    node[midway,above] {$x$};

\node[blue!70!black] at (4.5,3.35)
    {Electron displacement};

% Electric field
\draw[->,very thick,orange!90!black]
    (2.0,-0.45) -- (4.0,-0.45)
    node[midway,below] {$\mathbf E$};

% Restoring force
\draw[->,very thick,green!50!black]
    (4.8,-0.45) -- (4.0,-0.45)
    node[midway,below] {$\mathbf F=-e\mathbf E$};

\node[red!70!black,below] at (0.5,-1.05)
    {Stationary ions};
\node[blue!70!black,below] at (5.8,-1.05)
    {Displaced electrons};

\end{tikzpicture}
\caption{A displacement of the electron fluid relative to the
stationary ions produces oppositely charged surface layers.
The resulting electric field points from the positive sheet
toward the negative sheet, exerting a restoring force on
the electrons. The displacement is exaggerated for clarity.}
\label{fig:plasma_oscillation}
\end{figure}

By Gauss's law, $\nabla\cdot\mathbf E=\rho/\epsz$.
Equivalently, a Gaussian pillbox across a charged sheet gives
$\Delta E_\perp=\sigma_s/\epsz$.
Combining the fields of the two oppositely charged sheets gives
an electric field inside the slab of magnitude
\begin{equation}
E=\frac{\sigma_s}{\epsz}
=\frac{n_e e x}{\epsz}.
\end{equation}

For a positive displacement $x$, the field points to the right,
so the force $-eE$ on an electron points to the left:
it is a restoring force. Newton's law becomes
\begin{equation}
m_e\dv[2]{x}{t}
=-eE
=-\frac{n_e e^2}{\epsz}x.
\end{equation}

Thus
\begin{equation}
\boxed{
\dv[2]{x}{t}+\omega_p^2x=0,
\qquad
\omega_p^2=\frac{n_e e^2}{\epsz m_e}.
}
\end{equation}

The plasma frequency is the natural oscillation frequency of
this collective restoring motion.

\paragraph{Electron oscillation versus thermal motion.}

The electron fluid oscillates harmonically:
\begin{equation}
x(t)=x_0\cos(\omega_p t).
\end{equation}
Its maximum coherent oscillation velocity is therefore
\begin{equation}
\boxed{v_{\rm osc}=\omega_p x_0.}
\end{equation}

By the equipartition theorem, each quadratic degree of freedom
contributes $\frac{1}{2}k_B T_e$ to the average kinetic energy.
For one-dimensional thermal motion,
\begin{equation}
\boxed{
\frac{1}{2}m_e\langle v_x^2\rangle
=\frac{1}{2}k_B T_e
\quad\Longrightarrow\quad
v_{\rm th}=\sqrt{\frac{k_B T_e}{m_e}}.
}
\end{equation}


For example, at $T_e=10^4\,\mathrm{K}$,
\begin{equation}
v_{\rm th}\simeq 3.9\times10^5\,
\mathrm{m\,s^{-1}}
\simeq 390\,\mathrm{km\,s^{-1}}.
\end{equation}

The relative importance of the two velocities is
\begin{equation}
\boxed{
\frac{v_{\rm osc}}{v_{\rm th}}
=
\frac{\omega_p x_0}{\sqrt{k_B T_e/m_e}}.
}
\end{equation}

Unlike the thermal velocity, the oscillation velocity depends
on the amplitude of the disturbance. For sufficiently small
perturbations, $v_{\rm osc}$ can be much smaller than
$v_{\rm th}$.

Importantly, the \textbf{cold-plasma approximation does not
require $v_{\rm osc}>v_{\rm th}$}. Rather, it assumes that
thermal pressure and velocity-spread effects are negligible
in the wave dynamics.

For longitudinal waves, this typically requires
$k v_{\rm th}\ll\omega$, so that electrons do not move
appreciably across a wavelength during one oscillation.
For transverse electromagnetic waves, thermal corrections
to the usual cold, unmagnetized plasma dispersion relation
are likewise negligible in the appropriate long-wavelength,
nonrelativistic regime.

Thus, a plasma may have substantial electron thermal
motion and still be accurately described by the
cold-plasma approximation.


\subsection{Dispersion in cold plasma}
Now consider a \emph{collisionless, unmagnetized, cold plasma}: a neutral equilibrium with electron density $n_e$, stationary ions, and negligible electron thermal pressure. ``Cold'' means we neglect thermal motions in the linearized electron momentum equation; it does not mean the electrons are at absolute zero. Let $e>0$ be the magnitude of the electron charge, and $m_e$ its mass.

\paragraph{Step 1: Electron response to an electric field.}
We look for a single-frequency wave, so we write $\mathbf E\propto e^{-i\omega t}$. Why? In a linear, time-independent medium, sinusoidal Fourier components evolve independently; complex exponentials make time differentiation especially simple: $\partial_t\to-i\omega$. We take the real part at the end. This is an ansatz for a monochromatic component, not an assumption that every physical pulse has only one frequency.

For small-amplitude waves, the electron velocity $\mathbf v$ obeys
\begin{equation}
m_e\pdv{\mathbf v}{t}=-e\mathbf E.
\end{equation}
The magnetic Lorentz force $-e\mathbf v\times\mathbf B$ is second order in wave amplitude because both $\mathbf v$ and the wave magnetic field are first-order perturbations. For harmonic fields proportional to $e^{-i\omega t}$,
\begin{equation}
-i\omega m_e\mathbf v=-e\mathbf E
\quad\Longrightarrow\quad
\boxed{\mathbf v=-i\frac{e}{m_e\omega}\mathbf E.}
\end{equation}
Recall that $\mathbf J$ is the \textbf{electric current per unit area}, with SI units $\mathrm{A\,m^{-2}}$. For electrons with charge $-e$, number density $n_e$, and velocity $\mathbf v$, the current density is $\boxed{\mathbf J=-n_e e\mathbf v}$. Hence
\begin{equation}
\boxed{\mathbf J=i\frac{n_e e^2}{m_e\omega}\mathbf E.}
\end{equation}
Unlike an ordinary Ohmic conductor, the collisionless plasma has a frequency-dependent, imaginary effective conductivity.

\paragraph{Step 2: Derive the wave equation.}
Use Maxwell's equations with vacuum $\epsz,\muz$ and the plasma current:
\begin{align}
\nabla\times\mathbf E&=-\pdv{\mathbf B}{t},\\
\nabla\times\mathbf B&=\muz\mathbf J+\muz\epsz\pdv{\mathbf E}{t}.
\end{align}
For a transverse plane wave ($\mathbf k\cdot\mathbf E=0$), curl Faraday's law and substitute Amp\`ere's law:
\begin{equation}
\nabla^2\mathbf E=\muz\pdv{\mathbf J}{t}+\muz\epsz\pdv[2]{\mathbf E}{t}.
\end{equation}
Using the harmonic current above,
\begin{equation}
\pdv{\mathbf J}{t}=(-i\omega)\left(i\frac{n_e e^2}{m_e\omega}\mathbf E\right)
=\frac{n_e e^2}{m_e}\mathbf E.
\end{equation}
Therefore
\begin{equation}
\boxed{\nabla^2\mathbf E=\frac{1}{c^2}\pdv[2]{\mathbf E}{t}+\muz\frac{n_e e^2}{m_e}\mathbf E.}
\end{equation}
Substituting $\mathbf E=\mathbf E_0e^{i(kz-\omega t)}$ gives
\begin{equation}
-k^2=-\frac{\omega^2}{c^2}+\muz\frac{n_e e^2}{m_e}.
\end{equation}
Multiply by $c^2=1/(\muz\epsz)$, and recognize $\omega_p^2=n_e e^2/(\epsz m_e)$ from the preceding section:
\begin{equation}
\boxed{\omega^2=c^2k^2+\omega_p^2},\qquad
\boxed{\omega_p\equiv\sqrt{\frac{n_e e^2}{\epsz m_e}}}.
\end{equation}
Here $\omega_p$ is the electron plasma frequency derived above. Its ordinary-frequency form is $\nu_p=\omega_p/(2\pi)$.


\paragraph{Step 3: Index of refraction, phase and group velocities.}

The \textbf{index of refraction} $n$ is defined as the ratio of the speed of light in vacuum to the phase velocity of a wave in a medium:
\begin{equation}
\boxed{n\equiv\frac{c}{v_\mathrm{ph}}=\frac{ck}{\omega}.}
\end{equation}
Using the plasma dispersion relation $\omega^2=c^2k^2+\omega_p^2$, we obtain
\begin{equation}
\boxed{n=\sqrt{1-\frac{\omega_p^2}{\omega^2}}.}
\end{equation}
Unlike most ordinary materials, a cold plasma has $n<1$ for propagating electromagnetic waves.

From $k=\sqrt{\omega^2-\omega_p^2}/c$, the phase velocity is
\begin{equation}
\boxed{v_\mathrm{ph}=\frac{\omega}{k}
=\frac{c}{n}
=\frac{c}{\sqrt{1-\omega_p^2/\omega^2}}.}
\end{equation}
Differentiate $\omega=\sqrt{c^2k^2+\omega_p^2}$ to obtain the group velocity:
\begin{equation}
\boxed{v_g=\dv{\omega}{k}=\frac{c^2k}{\omega}
=cn=c\sqrt{1-\frac{\omega_p^2}{\omega^2}}.}
\end{equation}
Thus $v_\mathrm{ph}v_g=c^2$. Phase velocity can exceed $c$, but the group velocity of propagating waves is less than $c$. For $\omega<\omega_p$, $k$ is imaginary and transverse electromagnetic waves are evanescent rather than freely propagating, leading to reflection at a plasma boundary.


\textbf{Plasma reflection.}
Recall the cold-plasma dispersion relation,
\begin{equation}
\omega^2=c^2k^2+\omega_p^2.
\end{equation}
When $\omega<\omega_p$, the wavenumber $k$ becomes imaginary.
Electromagnetic waves can no longer propagate freely through the plasma;
instead, they become evanescent and are reflected at the plasma boundary
(in the ideal lossless case).

This is why the Earth's ionosphere can reflect radio waves back toward
the ground, particularly at frequencies below roughly $10\,\mathrm{MHz}$


\subsection{Dispersion measure}
A radio pulse contains many frequencies. In a dilute plasma, lower frequencies travel more slowly because $v_g$ depends on $\omega$. Suppose $\omega\gg\omega_p$ everywhere along the path. Expand $\sqrt{1-x}\simeq1-x/2$ for $x\ll1$:
\begin{equation}
v_g=c\sqrt{1-\frac{\omega_p^2}{\omega^2}}
\simeq c\left(1-\frac{\omega_p^2}{2\omega^2}\right).
\end{equation}
\textbf{Important:} We need the \emph{reciprocal} of the group velocity to calculate travel time, not the group velocity itself. Using $1/(1-a)\simeq1+a$ for $|a|\ll1$, the sign changes:
\begin{equation}
\frac{1}{v_g}=\frac{1}{c}\left(1-\frac{\omega_p^2}{\omega^2}\right)^{-1/2}
\simeq\frac{1}{c}\left(1+\frac{\omega_p^2}{2\omega^2}\right).
\end{equation}
For an element of path length $d\ell$, the extra travel time relative to vacuum is
\begin{equation}
dt_\mathrm{extra}=\left(\frac{1}{v_g}-\frac{1}{c}\right)d\ell
\simeq\frac{\omega_p^2}{2c\omega^2}\,d\ell.
\end{equation}
Substitute $\omega_p^2=n_e e^2/(\epsz m_e)$ and $\omega=2\pi\nu$:
\begin{equation}
\boxed{\Delta t(\nu)=\frac{e^2}{8\pi^2\epsz m_e c\nu^2}\int n_e(\ell)\,d\ell.}
\end{equation}
Define the \textbf{dispersion measure} (DM) as the electron column density along the line of sight:
\begin{equation}
\boxed{\mathrm{DM}\equiv\int_0^D n_e(\ell)\,d\ell.}
\end{equation}
Then
\begin{equation}
\boxed{\Delta t(\nu)=\frac{e^2}{8\pi^2\epsz m_e c}\frac{\mathrm{DM}}{\nu^2}.}
\end{equation}
For two observing frequencies $\nu_\mathrm{low}<\nu_\mathrm{high}$,
\begin{equation}
\boxed{t(\nu_\mathrm{low})-t(\nu_\mathrm{high})
=K_\mathrm{DM}\,\mathrm{DM}\left(\nu_\mathrm{low}^{-2}-\nu_\mathrm{high}^{-2}\right).}
\end{equation}
A convenient radio-astronomy convention is
\begin{equation}
\boxed{\Delta t\,[\mathrm{ms}]\simeq4.1488\,\mathrm{DM}\,[\mathrm{pc\,cm^{-3}}]
\left[\left(\frac{\nu_\mathrm{low}}{\mathrm{GHz}}\right)^{-2}
-\left(\frac{\nu_\mathrm{high}}{\mathrm{GHz}}\right)^{-2}\right].}
\end{equation}
This is the familiar $\nu^{-2}$ dispersive sweep seen in pulsars and fast radio bursts. The derivation assumes a cold, weakly dispersive, unmagnetized plasma and neglects cosmological redshift; for cosmological sources, the observed contribution of a plasma element at redshift $z$ is weighted by $1/(1+z)$ when written in terms of its proper electron column.



For an element of path length $d\ell$, the extra travel time relative to vacuum is
\begin{equation}
dt_\mathrm{extra}=\left(\frac{1}{v_g}-\frac{1}{c}\right)d\ell
\simeq\frac{\omega_p^2}{2c\omega^2}\,d\ell.
\end{equation}
Substitute $\omega_p^2=n_e e^2/(\epsz m_e)$ and $\omega=2\pi\nu$:
\begin{equation}
\boxed{\Delta t(\nu)=\frac{e^2}{8\pi^2\epsz m_e c\nu^2}\int n_e(\ell)\,d\ell.}
\end{equation}
Define the \textbf{dispersion measure} (DM) as the electron column density along the line of sight:
\begin{equation}
\boxed{\mathrm{DM}\equiv\int_0^D n_e(\ell)\,d\ell.}
\end{equation}
Then
\begin{equation}
\boxed{\Delta t(\nu)=\frac{e^2}{8\pi^2\epsz m_e c}\frac{\mathrm{DM}}{\nu^2}.}
\end{equation}
For two observing frequencies $\nu_\mathrm{low}<\nu_\mathrm{high}$,
\begin{equation}
\boxed{t(\nu_\mathrm{low})-t(\nu_\mathrm{high})
=K_\mathrm{DM}\,\mathrm{DM}\left(\nu_\mathrm{low}^{-2}-\nu_\mathrm{high}^{-2}\right).}
\end{equation}
A convenient radio-astronomy convention is
\begin{equation}
\boxed{\Delta t\,[\mathrm{ms}]\simeq4.1488\,\mathrm{DM}\,[\mathrm{pc\,cm^{-3}}]
\left[\left(\frac{\nu_\mathrm{low}}{\mathrm{GHz}}\right)^{-2}
-\left(\frac{\nu_\mathrm{high}}{\mathrm{GHz}}\right)^{-2}\right].}
\end{equation}

This is the familiar $\nu^{-2}$ dispersive sweep seen in pulsars and fast radio bursts.

\paragraph{Cosmological redshift correction.}
For a plasma element at redshift $z$, the local wave frequency is
$\nu_\mathrm{local}=(1+z)\nu_\mathrm{obs}$, while the observed time interval
is stretched by a factor $(1+z)$. Together, these give a net factor
of $1/(1+z)$.

Thus, for a cosmological source,
\begin{equation}
\boxed{
\Delta t_\mathrm{obs}(\nu_\mathrm{obs})
=\frac{K_\mathrm{DM}}{\nu_\mathrm{obs}^2}
\int_0^{z_s}\frac{n_e(z)}{1+z}\,d\ell_\mathrm{proper}.
}
\end{equation}
We can therefore define an observed dispersion measure,
\begin{equation}
\boxed{
\mathrm{DM}_\mathrm{obs}
=\int_0^{z_s}\frac{n_e(z)}{1+z}\,d\ell_\mathrm{proper}.
}
\end{equation}




